% =========================================================
%  sesi/05/sesi.tex
%  Sesi 5 — Sprite, artefak, peti, dan pintu
% =========================================================

\sesi{05}{Sprite, Artefak, Peti, dan Pintu}{1.4}{140 menit}

Sesi ini adalah pertemuan minggu ke-5. Pada sesi sebelumnya pemain sudah
jatuh karena \istilah{gravitasi}, melompat dengan \istilah{jump impulse}, dan
memantul di dinding dengan \istilah{koefisien restitusi} 0,8. Namun seluruh isi
gudang masih berupa daftar persegi panjang: lantai, rak, dan pemain saja.
Belum ada yang dapat diambil, belum ada yang berbahaya, dan belum ada pintu
keluar.

Pada sesi ini gudang diisi. Enam artefak tersebar di atas rak, tiga peti
menghalangi jalan di lantai, dan sebuah pintu menunggu di ujung kanan. Pemain
mengumpulkan artefak untuk menambah skor sepuluh poin, menghindari peti yang
mengurangi satu nyawa, lalu keluar setelah seluruh artefak terkumpul. Nyawa
awal tiga; bila habis, permainan berakhir.

Untuk mengelola benda yang jumlahnya bertambah itu, Anda memakai
\istilah{sprite} dan \istilah{group} dari Pygame, serta memeriksa persentuhan
dengan \istilah{AABB} berupa \code{Rect}. Urutan \istilah{integrasi Euler} dan
penyelesaian tabrakan per sumbu yang Anda pelajari pada sesi 4 tetap dipakai;
yang berubah adalah bentuk data yang diperiksa.

% ---------------------------------------------------------
\bagian{Tujuan Sesi}

\begin{tujuanbox}
Setelah menuntaskan sesi ini, Anda dapat menunjukkan:
\begin{enumerate}[leftmargin=1.5em, itemsep=0.2em]
  \item lantai, rak, artefak, peti, dan pintu sebagai kelas turunan
        \code{pygame.sprite.Sprite} yang dikelola \code{Group};
  \item artefak yang hilang ketika tersentuh dan menambah skor sepuluh poin
        pada HUD;
  \item peti yang mengurangi satu nyawa dari tiga nyawa awal dan mengembalikan
        pemain ke titik awal tanpa kehilangan nyawa berulang;
  \item pintu yang tetap tertutup sampai artefak terakhir terkumpul, lalu
        membuka dan mengakhiri level ketika disentuh;
  \item penjelasan mengapa tabrakan diperiksa dengan \code{colliderect} dan
        \code{spritecollide}, serta mengapa hasilnya harus diselesaikan per
        sumbu.
\end{enumerate}
\end{tujuanbox}

% ---------------------------------------------------------
\bagian{Prasyarat}

\begin{prasyaratbox}
\begin{itemize}[leftmargin=1.4em, itemsep=0.25em]
  \item Anda sudah menjalankan checkpoint sesi 4: pemain berjalan, jatuh,
        mendarat di rak, melompat, dan memantul di dinding.
  \item Anda memahami perbedaan \code{position} bertipe pecahan dan
        \code{rect} bertipe bilangan bulat.
  \item Anda pernah menulis kelas Python sederhana dengan \code{__init__} dan
        \code{self}.
  \item Anda dapat menjelaskan mengapa pemeriksaan tabrakan dilakukan sesudah
        posisi diperbarui, bukan sebelumnya.
\end{itemize}
\end{prasyaratbox}

% ---------------------------------------------------------
\bagian{Peta Kemampuan}

\petasesi
  {Pemain dan fisika dari sesi 4; lantai dan rak masih berupa daftar
   \code{Rect}; belum ada benda yang dapat diambil.}
  {Kelas \code{Sprite} untuk seluruh isi level, wadah \code{Group},
   \code{spritecollide}, \istilah{AABB}, skor, nyawa, pintu, dan HUD.}
  {NPC yang bergerak dan mengambil keputusan sendiri pada sesi 6.}

% ---------------------------------------------------------
\bagian{Konsep yang Dipakai}

\subsection*{Sprite: gambar dan persegi pembatas}

Sebuah \istilah{sprite} adalah benda permainan yang menyatukan gambar dan
letaknya. Pygame menyediakan kelas dasar \code{pygame.sprite.Sprite}; kelas
turunan Anda wajib memanggil \code{super().__init__()} dan mengisi dua nama
yang sudah disepakati: \code{self.image} berupa \code{Surface}, dan
\code{self.rect} berupa \code{Rect}. Nama itu tidak boleh diganti, sebab
perkakas Pygame mencarinya ketika menggambar dan memeriksa tabrakan.

\subsection*{Group: satu perintah untuk banyak benda}

\istilah{Group} adalah wadah yang menyimpan sekumpulan sprite. Perintah
\code{group.draw(screen)} menggambar seluruh anggotanya,
\code{group.update(...)} memanggil pembaruan setiap anggota, dan
\code{for benda in group} menelusurinya seperti daftar biasa. Satu sprite
boleh berada di beberapa group sekaligus, karena group hanya menyimpan
rujukan, bukan salinan. Pada sesi ini ada tiga group: \code{solids},
\code{artefacts}, dan \code{crates}.

\subsection*{AABB dan arti persentuhan}

\istilah{AABB} (\emph{axis-aligned bounding box}) adalah persegi pembatas yang
sisinya sejajar sumbu layar. Pemain berukuran 36 kali 36 piksel, sehingga
persegi pembatasnya berimpit dengan gambarnya. Pemeriksaannya memakai
\code{Rect}, sehingga murah dan cukup untuk benda yang bentuknya persegi.
Perlu diingat bahwa
\code{colliderect} hanya bernilai benar bila kedua persegi beririsan dengan
luas lebih dari nol; dua persegi yang hanya bersentuhan pada tepinya dinilai
tidak bertabrakan. Sifat itu berguna agar benda yang sudah ditempelkan pada
tepi tidak dianggap bertabrakan lagi.

\subsection*{Mendeteksi berbeda dari menyelesaikan}

\code{colliderect} menjawab ``apakah kedua persegi bersinggungan?''.
\code{spritecollide(benda, group, dokill)} menjawab pertanyaan yang sama untuk
satu benda terhadap sekumpulan benda, dan mengembalikan daftar. Argumen
\code{dokill} menentukan apakah benda yang tersentuh langsung dikeluarkan dari
group. Untuk artefak nilainya benar; untuk peti nilainya salah, sebab peti
tetap berada di tempatnya.

Pertanyaan berikutnya---dari sisi mana tabrakan terjadi dan apa yang harus
dilakukan---tidak dijawab oleh pustaka. Itulah \istilah{collision} yang harus
Anda selesaikan sendiri: pindahkan pada satu sumbu, periksa, perbaiki posisi
pada sumbu itu, lalu lanjut ke sumbu berikutnya.

% ---------------------------------------------------------
\bagian{Tutorial}

Kerjakan langkah berikut berurutan. Pada langkah pertama, Anda memindahkan
benda yang sebelumnya berupa \code{Rect} menjadi kelas \code{Sprite}; jangan
mencampur kedua cara dalam satu group.

\begin{langkah}
  \item \textbf{Ubah isi level menjadi sprite.} Di
        \file{gudang_artefak/entities.py}, tulis kelas \code{Solid},
        \code{Artefact}, \code{Crate}, dan \code{Door}. Masing-masing
        memanggil \code{super().__init__()}, membuat \code{self.image}, lalu
        menetapkan \code{self.rect}. Artefak digambar sebagai belah ketupat
        berukuran 22 piksel, peti sebagai persegi 34 piksel bergaris silang,
        dan pintu sebagai persegi panjang 46 kali 78 piksel.

  \item \textbf{Berikan keadaan pada pintu.} Tambahkan \code{self.open =
        False} dan metode \code{set_open(value)} pada \code{Door}. Metode itu
        menggambar ulang pintu hanya ketika keadaannya berubah, sehingga pintu
        tertutup berwarna abu-abu dan pintu terbuka berwarna hijau.

  \item \textbf{Susun level sebagai kamus.} Di \file{levels.py}, tulis
        \code{build_rak()} yang mengembalikan kamus berisi \code{solids},
        \code{artefacts}, \code{crates}, \code{door}, \code{spawn}, dan
        \code{name}. Isi setiap group dari daftar titik yang sudah ada:
        enam artefak, tiga peti, dan satu pintu.

  \item \textbf{Siapkan keadaan permainan.} Di \code{Game}, tulis
        \code{reset()} yang menetapkan \code{score = 0},
        \code{lives = S.START_LIVES} (tiga), \code{mode = "rak"}, dan
        \code{artefacts_left} sejumlah isi \code{ARTEFACT_LAYOUT}. Panggil
        \code{reset()} dari \code{__init__} agar permainan selalu mulai dari
        keadaan yang sama.

  \item \textbf{Tambahkan \code{respawn}.} Pada \code{Player}, tulis metode
        yang mengembalikan posisi dan \code{rect} ke titik awal, mengosongkan
        kecepatan, dan memberi \code{respawn_timer} satu detik. Waktu kebal itu
        mencegah peti yang sama mengurangi nyawa berkali-kali dalam sekejap.

  \item \textbf{Perbarui level Rak.} Pindahkan pemain dengan
        \code{update_side(delta_seconds, direction, self.solids)}. Karena
        \code{solids} kini berupa \code{Group}, pemeriksaan di dalam
        \code{Player} memakai \code{solid.rect}, bukan \code{solid} langsung.

  \item \textbf{Kumpulkan artefak.} Panggil
        \code{spritecollide(player, artefacts, True)}. Telusuri daftar hasilnya;
        setiap anggota menambah sepuluh poin dan mengurangi
        \code{artefacts_left}. Sesudah itu, panggil
        \code{door.set_open(artefacts_left == 0)}.

  \item \textbf{Periksa peti.} Bila \code{respawn_timer} sudah habis, panggil
        \code{spritecollide(player, crates, False)}. Bila hasilnya tidak
        kosong, kurangi nyawa. Bila nyawa masih tersisa, kembalikan pemain ke
        titik awal; bila habis, ubah mode menjadi kalah.

  \item \textbf{Buka pintu dan gambar HUD.} Periksa
        \code{door.open and player.rect.colliderect(door.rect)} untuk
        mengakhiri level. Gambar ketiga group dengan \code{draw}, lalu tampilkan
        nama level, skor, dan nyawa pada sudut kiri atas.
\end{langkah}

\begin{ujibox}
Setelah Langkah 9, periksa lima hal berikut.
\begin{itemize}[leftmargin=1.4em, itemsep=0.2em]
  \item Seluruh enam artefak terlihat di atas rak, dan menyentuhnya membuat
        artefak itu hilang serta skor bertambah sepuluh.
  \item Setelah artefak keenam, skor bernilai 60 dan pintu berubah menjadi
        hijau. Sebelum itu, pintu tetap abu-abu dan tidak dapat dibuka.
  \item Menyentuh peti mengurangi nyawa dari tiga menjadi dua, dan pemain
        kembali ke titik awal.
  \item Berdiri menempel pada peti tidak mengurangi nyawa berkali-kali; waktu
        kebal satu detik bekerja.
  \item Menyentuh pintu yang sudah terbuka menampilkan pesan
        \code{Level Rak selesai.}, dan menekan R memulai ulang permainan.
\end{itemize}
\end{ujibox}

% ---------------------------------------------------------
\bagian{Kode Sesi Ini}

Potongan berikut adalah bagian yang berubah pada checkpoint sesi ini.
Perhatikan bahwa \code{Player} sudah berupa sprite sejak sesi 3; yang baru
adalah seluruh isi level mengikuti bentuk yang sama, lalu diperiksa dengan
perkakas group.

\tampilkodebaris[Python]{code/sesi-05/gudang_artefak/entities.py}{107}{117}
  {entities.py: artefak sebagai sprite yang dapat diambil}

\tampilkodebaris[Python]{code/sesi-05/gudang_artefak/entities.py}{120}{128}
  {entities.py: peti sebagai sprite yang tetap di tempatnya}

\tampilkodebaris[Python]{code/sesi-05/gudang_artefak/entities.py}{131}{154}
  {entities.py: pintu dengan keadaan terbuka dan tertutup}

\tampilkodebaris[Python]{code/sesi-05/gudang_artefak/entities.py}{88}{93}
  {entities.py: respawn dan waktu kebal sesudah kehilangan nyawa}

\tampilkodebaris[Python]{code/sesi-05/gudang_artefak/levels.py}{32}{49}
  {levels.py: menyusun seluruh isi level menjadi group dan kamus}

\tampilkodebaris[Python]{code/sesi-05/gudang_artefak/game.py}{93}{118}
  {game.py: artefak, peti, dan pintu diperiksa setiap bingkai}

\tampilkodebaris[Python]{code/sesi-05/gudang_artefak/game.py}{120}{126}
  {game.py: kehilangan nyawa dan keputusan mengakhiri permainan}

\tampilkodebaris[Python]{code/sesi-05/gudang_artefak/game.py}{150}{155}
  {game.py: HUD nama level, skor, dan nyawa}

\checkpoint{05}{Enam artefak dapat dikumpulkan dengan nilai sepuluh poin
masing-masing, peti mengurangi satu dari tiga nyawa dan mengembalikan pemain,
pintu terbuka setelah artefak terakhir, dan HUD menampilkan skor serta nyawa.}

% ---------------------------------------------------------
\bagian{Kesalahan Umum}

\begin{galatbox}
\begin{itemize}[leftmargin=1.4em, itemsep=0.35em]
  \item \textbf{Benda tidak tergambar atau tidak pernah bertabrakan.}
        Gejala: group tampak kosong, atau artefak tidak pernah hilang. Periksa
        tiga hal pada kelas sprite: \code{super().__init__()} dipanggil, nama
        \code{image} dan \code{rect} tepat, dan \code{rect} benar-benar
        berupa \code{Rect}.
  \item \textbf{\code{TypeError} saat memanggil \code{spritecollide}.}
        Gejala: program berhenti dengan pesan bahwa argumen bukan group.
        Periksa bahwa argumen kedua adalah \code{pygame.sprite.Group}, bukan
        daftar biasa. Ubah daftar menjadi group dengan
        \code{pygame.sprite.Group(daftar)}.
  \item \textbf{Skor bertambah setiap bingkai selama artefak disentuh.}
        Gejala: skor melonjak tanpa artefak menghilang. Periksa argumen ketiga
        \code{spritecollide}; untuk artefak nilainya harus benar agar benda
        dikeluarkan dari group.
  \item \textbf{Nyawa habis dalam sekejap ketika menyentuh peti.}
        Gejala: tiga nyawa hilang pada satu sentuhan. Periksa
        \code{respawn_timer}; tanpa waktu kebal, peti yang sama diperiksa dan
        mengenai pemain pada setiap bingkai.
  \item \textbf{Pintu sudah terbuka sejak awal.}
        Gejala: pemain dapat menang tanpa mengumpulkan artefak. Periksa
        \code{artefacts_left} pada \code{reset()} dan pastikan
        \code{door.set_open(False)} dipanggil ketika level dimuat.
  \item \textbf{Pemain tetap menempel pada peti sesudah respawn.}
        Gejala: pemain kembali ke titik awal tetapi langsung meluncur atau
        terjebak. Periksa bahwa \code{respawn} mengosongkan
        \code{velocity} dan menyelaraskan \code{rect} dengan posisi baru.
\end{itemize}
\end{galatbox}

% ---------------------------------------------------------
\bagian{Lembar Kerja}

Isi tabel berikut pada catatan praktikum Anda. Baris keempat diisi dengan
mencoba membuka pintu sebelum seluruh artefak terkumpul.

{\small
\renewcommand{\arraystretch}{1.4}
\begin{tabularx}{\linewidth}{@{}C{0.8cm} L{5.4cm} Y@{}}
\toprule
\textbf{No} & \textbf{Hasil yang dicatat} & \textbf{Keterangan Anda} \\
\midrule
1 & Jumlah artefak dan skor maksimum yang mungkin & \\[0.55cm]
2 & Nyawa awal dan nyawa setelah satu peti & \\[0.55cm]
3 & Posisi pemain sesudah respawn & \\[0.55cm]
4 & Keadaan pintu sebelum dan sesudah artefak keenam & \\[0.55cm]
5 & Satu kendala yang Anda hadapi dan cara menyelesaikannya & \\[0.9cm]
\bottomrule
\end{tabularx}
}

Catat pula berapa bingkai yang dibutuhkan pemain untuk menyentuh peti yang
sama dua kali. Angka itu menunjukkan apakah waktu kebal satu detik cukup
panjang pada laju bingkai yang Anda pakai.

% ---------------------------------------------------------
\bagian{Checklist Sesi}

\begin{checklistbox}
\begin{checklist}
  \item Saya dapat menunjukkan seluruh isi level sebagai sprite yang dikelola
        group.
  \item Saya dapat menjelaskan dua nama yang wajib ada pada setiap sprite dan
        akibatnya bila salah nama.
  \item Saya dapat memilih nilai \code{dokill} yang tepat untuk artefak dan
        untuk peti.
  \item Saya dapat menunjukkan artefak yang menambah sepuluh poin dan HUD yang
        ikut berubah.
  \item Saya dapat menunjukkan peti yang mengurangi satu nyawa dan
        mengembalikan pemain tanpa pengurangan berulang.
  \item Saya dapat menjelaskan mengapa pintu baru terbuka setelah
        \code{artefacts_left} bernilai nol.
\end{checklist}
\end{checklistbox}

% ---------------------------------------------------------
\bagian{Pengait ke Sesi Berikutnya}

Gudang sekarang dapat dimainkan, tetapi seluruh benda diam di tempatnya dan
hanya pemain yang bergerak. Pada sesi 6 Anda menambahkan penjaga: sprite yang
melihat posisi pemain, menghitung arah menuju pemain, lalu bergerak sendiri.
Perkakas yang baru Anda pakai---kelas \code{Sprite}, \code{Group}, dan
pemeriksaan \code{Rect}---tetap dipakai, sehingga penjaga cukup ditambahkan
sebagai anggota group baru tanpa mengubah cara kerja level.
